\head{第三卷第二部分：出版信息与书末索引}
\source{EGAIII.2.pdf，封面及书末；原印刷第81--83页（通卷页码213--215）}
\noindent\textbf{书名。}《代数几何基础》第III章：《凝聚层的上同调研究》，第二部分。A. Grothendieck 著，与 J. Dieudonné 合作撰写。法国高等科学研究所，1963年；\textit{Publications mathématiques de l'I.H.É.S.}，第17卷，第5--91页。初版法定送存：1963年第四季度。各国版权均保留。©1963，Institut des Hautes Études Scientifiques。

\noindent\textbf{底本的数字化版权说明（译文）。}©\textit{Publications mathématiques de l'I.H.É.S.}，1963，保留所有权利。访问该刊档案意味着接受其一般使用条件。任何商业使用或系统性印刷均构成违法行为。本文件的任何复制件或印刷件均须包含此版权声明。本文由“古数学文献数字化”计划数字化。

\NewNote{以下保留本分册原有索引的组织和指向；号码为原著条目号，不是本合订本页码。上传底本在正文之后依次为记号索引、术语索引、原目录、空白页和《勘误与增补（清单二）》，没有另列参考文献续表，故不补造一份续表。正文末的“待续”是原著历史标记。}

\section*{记号索引}
\addcontentsline{toc}{section}{第三卷第二部分记号索引}
\begin{longtable}{@{}p{.77\linewidth}p{.18\linewidth}@{}}
\toprule 记号 & 原著条目\\\midrule\endhead
$\operatorname{Tor}_n^A(P_\bullet,Q_\bullet)$ & 6.3.1\\
$\mathfrak{Tor}_n^S(\PP_\bullet,\mathscr Q_\bullet)$、$\mathfrak{Tor}_n(\PP_\bullet,\mathscr Q_\bullet)$ & 6.4.1；6.5.1\\
$\lTor_n^S(\F,\G)$ & 6.5.1\\
$\mathfrak{Tor}_q^S(\PP_\bullet^{(1)},\PP_\bullet^{(2)},\ldots,\PP_\bullet^{(m)})$ & 6.5.15\\
$H^{-n}(\mathfrak U^{(i)},\PP_\bullet^{(i)})$、$H^{-n}(X^{(i)},\PP_\bullet^{(i)})$ & 6.6.1\\
$\operatorname{Tor}_n^A(L_{\bullet\bullet}^{(1)},L_{\bullet\bullet}^{(2)})$；$\operatorname{Tor}_n^S(\mathfrak U^{(1)},\mathfrak U^{(2)};\PP_\bullet^{(1)},\PP_\bullet^{(2)})$ & 6.6.2\\
$\lTor_n^S(\mathfrak U^{(1)},\mathfrak U^{(2)};\F^{(1)},\F^{(2)})$；$\lTor_n^S(\mathfrak U^{(1)},\mathfrak U^{(2)};\PP_\bullet^{(1)},\PP_\bullet^{(2)})$ & 6.6.2\\
$\mathfrak{Tor}_n^S(f_1,f_2;\PP_\bullet^{(1)},\PP_\bullet^{(2)})$ & 6.6.8；6.7.1；6.7.3\\
$\mathfrak{Tor}_n^S((f_i)_{i\in I};(\PP_\bullet^{(i)})_{i\in I})$；$\mathfrak{Tor}_n^S(f_1,\ldots,f_m;\PP_\bullet^{(1)},\ldots,\PP_\bullet^{(m)})$ & 6.7.11\\
$\mathbf{Ab}$、$\mathbf{Ab}_A$ & 7.1.1\\
$T_{(B)}$、$T^{(B)}$、$T\otimes_A B$（$T$ 为从 $\mathbf{Ab}_A$ 到 $\mathbf{Ab}$ 的函子） & 7.1.3\\
$t_M$ & 7.2.2\\
$T_{\mathfrak p}$ & 7.1.4\\
$\TT_\bullet^{Y'}(\M')$、$\TT_\bullet^U(\M|U)$ & 7.7.2\\
$\operatorname{EP}(f,\PP_\bullet;y)$、$\operatorname{EP}(Y,\PP_\bullet;y)$、$\operatorname{EP}(\PP_\bullet;y)$ & 7.9.5\\
$\operatorname{PH}(f,\PP_\bullet;y)$、$\operatorname{PH}(\PP_\bullet;y)$ & 7.9.12\\
\bottomrule
\end{longtable}
\NewNote{原记号索引以不同字形区分超 Tor 和逐项 Tor。本译稿正文采用 $\bTor$ 与带 term 标识的记号来避免字体难辨，并在定义处说明。此表保留原索引的区分及其条目指向；Hilbert 多项式的命名在（7.9.11）末出现，但原索引指向（7.9.12），此处不静默改号。}

\section*{术语索引}
\addcontentsline{toc}{section}{第三卷第二部分术语索引}
\begin{longtable}{@{}p{.77\linewidth}p{.18\linewidth}@{}}
\toprule 术语 & 原著条目\\\midrule\endhead
模复形的 Euler--Poincaré 特征数 & 7.9.5\\
在 $y\in Y$ 处关于 $Y$ 在维数 $p$ 上上同调平坦；关于 $Y$ 在维数 $p$ 上上同调平坦；在 $y$ 处关于 $Y$ 上同调平坦；关于 $Y$ 上同调平坦（$Y$-平坦的 $\OO_X$-模复形） & 7.8.1\\
同伦等价的复形 & 6.1.4\\
上同调维数 $\leq n$ 的赋环空间 & 6.5.5\\
协变加性函子的标量扩张 & 7.1.3\\
上同调维数 $\leq n$ 的环层 & 6.5.5\\
Künneth 公式 & 6.7.8\\
在 $y\in Y$ 处关于 $Y$ 在维数 $p$ 上同调平坦；关于 $Y$ 在维数 $p$ 上同调平坦；在 $y$ 处关于 $Y$ 同调平坦；关于 $Y$ 同调平坦（$Y$-平坦的 $\OO_X$-模复形） & 7.8.1\\
同伦等价态射 & 6.1.4\\
两个 $A$-模复形的超 Tor & 6.3.1\\
两个模层复形的局部超 Tor & 6.4.1\\
两个模层复形的整体超 Tor & 6.6.2；6.6.8；6.7.3\\
模复形的 Hilbert 多项式 & 7.9.12\\
Künneth 谱序列 & 6.7.3\\
\bottomrule
\end{longtable}

\section*{原著目录}
\addcontentsline{toc}{section}{第三卷第二部分原著目录}
\NewNote{右栏页码是该分册原印刷页码，专供与底本核对；本合订本的实际页码见总目录。}
\begin{longtable}{@{}p{.86\linewidth}r@{}}
\toprule 内容 & 原页\\\midrule\endhead
第III章：凝聚层的上同调研究（续）&\\
§6．局部与整体 Tor 函子；Künneth 公式 &5\\
6.1．引言 &5\\
6.2．预概形上模复形的超上同调 &6\\
6.3．两个模复形的超 Tor &9\\
6.4．拟凝聚模层复形的局部超 Tor 函子：仿射概形情形 &14\\
6.5．拟凝聚模层复形的局部超 Tor 函子：一般情形 &16\\
6.6．拟凝聚模层复形的整体超 Tor 函子与 Künneth 谱序列：仿射底情形 &21\\
6.7．拟凝聚模层复形的整体超 Tor 函子与 Künneth 谱序列：一般情形 &25\\
6.8．整体超 Tor 的结合性谱序列 &32\\
6.9．整体超 Tor 的基变换谱序列 &34\\
6.10．某些上同调函子的局部结构 &39\\
§7．模的协变同调函子中的基变换研究 &43\\
7.1．$A$-模函子 &43\\
7.2．张量积函子的刻画 &44\\
7.3．模的同调函子的正合性判据 &48\\
7.4．函子 $H_\bullet(P_\bullet\otimes_A M)$ 的正合性判据 &53\\
7.5．诺特局部环的情形 &58\\
7.6．正合性的下降；半连续性定理与 Grauert 正合性判据 &60\\
7.7．对固有态射的应用：I．交换性质 &65\\
7.8．对固有态射的应用：II．上同调平坦性判据 &72\\
7.9．对固有态射的应用：III．Euler--Poincaré 特征数与 Hilbert 多项式的不变性 &76\\
\bottomrule
\end{longtable}
\begin{flushright}收稿日期：1962年12月15日。\end{flushright}
